% =====================================================================
%  Section 2.7 - PIN photodetectors  &  Section 2.8 - Optical splitters
%  To be \input{} inside the main report, under the Section 2 heading.
%  Requires: amsmath, siunitx, booktabs
% =====================================================================

% ---------------------------------------------------------------------
\subsection{PIN photodetectors}
\label{sec:2_7}
% ---------------------------------------------------------------------

\paragraph{Data.}
Quantum efficiency $\eta_\mathrm{ph} = 0.8$, optical frequency $f_c = \SI{193.5}{\tera\hertz}$,
incident (received) power $P_\mathrm{RX} = P_\mathrm{inc} = \SI{-20}{dBm}$, bitrate $R_b = \SI{50}{\mega\bit\per\second}$,
TIA noise standard deviation $\sigma_\mathrm{TIA} = \sigma_n = \SI{0.5}{\micro\ampere}$.

\noindent Constants: $h = \SI{6.626e-34}{\joule\second}$, $q = \SI{1.602e-19}{\coulomb}$.

\noindent First, convert the incident power to linear units:
\begin{equation}
  P_\mathrm{RX} = \SI{1}{\milli\watt}\cdot 10^{-20/10} = \SI{e-5}{\watt} = \SI{10}{\micro\watt}.
\end{equation}

\paragraph{1. Energy of one photon.}
A photon at frequency $f_c$ carries the energy
\begin{equation}
  E_\mathrm{ph} = h f_c = \SI{6.626e-34}{\joule\second}\times\SI{193.5e12}{\hertz}
  = \SI{1.282e-19}{\joule} \;\approx\; \SI{0.80}{\electronvolt}.
\end{equation}

\paragraph{2. Average number of incident photons per second.}
The photon flux is the optical power divided by the energy per photon:
\begin{equation}
  \bar N_\mathrm{ph} = \frac{P_\mathrm{RX}}{h f_c}
  = \frac{\SI{e-5}{\watt}}{\SI{1.282e-19}{\joule}}
  = \SI{7.80e13}{photons\per\second}.
\end{equation}

\paragraph{3. Responsivity and photocurrent.}
Each photon generates on average $\eta_\mathrm{ph}$ electrons, so the current is
$i = \eta_\mathrm{ph}\, q\, \bar N_\mathrm{ph} = \eta_\mathrm{ph}\, q\, P_\mathrm{RX}/(h f_c)$.
Writing $i_\mathrm{RX} = R\, P_\mathrm{RX}$, the responsivity is
\begin{equation}
  R = \frac{\eta_\mathrm{ph}\, q}{h f_c}
  = \frac{0.8\times\SI{1.602e-19}{\coulomb}}{\SI{1.282e-19}{\joule}}
  = \SI{0.9996}{\ampere\per\watt} \;\approx\; \SI{1}{\ampere\per\watt},
\end{equation}
and the photocurrent is
\begin{equation}
  i_\mathrm{RX} = R\, P_\mathrm{RX} = \SI{0.9996}{\ampere\per\watt}\times\SI{e-5}{\watt}
  \approx \SI{10}{\micro\ampere}.
\end{equation}

\paragraph{4. Average number of incident photons per bit.}
The bit duration is $T_b = 1/R_b = \SI{20}{\nano\second}$, therefore
\begin{equation}
  \bar N_\mathrm{ph/bit} = \frac{\bar N_\mathrm{ph}}{R_b} = \frac{P_\mathrm{RX}\,T_b}{h f_c}
  = \frac{\SI{7.80e13}{\per\second}}{\SI{50e6}{\per\second}}
  = \num{1.56e6}\ \text{photons/bit}.
\end{equation}

\paragraph{5. Electrical SNR of the PIN receiver.}
Assuming the receiver is limited by the TIA noise (thermal-noise-limited receiver),
the electrical SNR is the ratio between the signal power and the noise power,
both proportional to the square of the respective currents. The TIA noise is given directly through its
standard deviation, which is equivalent to $\sigma_\mathrm{TIA}^2 = \mathrm{IRND}^2\cdot B_\mathrm{eq}$:
\begin{equation}
\begin{aligned}
  \mathrm{SNR}^\mathrm{PIN+TIA}_\mathrm{TOT} &= \frac{i_\mathrm{RX}^2}{\sigma_\mathrm{TIA}^2}
  = \frac{(R\,P_\mathrm{RX})^2}{\sigma_\mathrm{TIA}^2}
  = \frac{(\SI{10}{\micro\ampere})^2}{(\SI{0.5}{\micro\ampere})^2}
  = \frac{\SI{e-10}{\ampere\squared}}{\SI{2.5e-13}{\ampere\squared}}
  \approx 400, \\
  \mathrm{SNR}^\mathrm{PIN+TIA}_\mathrm{TOT,dB} &= 10\log_{10}(400) \approx \SI{26.0}{dB}.
\end{aligned}
\end{equation}
\emph{Check of the assumption.} The PIN shot-noise variance over $B_\mathrm{eq}=\SI{50}{\mega\hertz}$ is
$\sigma_\mathrm{shot}^2 = 2qi_\mathrm{RX}B_\mathrm{eq} = \SI{1.6e-16}{\ampere\squared}$,
about $0.06\%$ of $\sigma_\mathrm{TIA}^2$: it is negligible, so the TIA-noise-only
model is justified.

\paragraph{APD receiver.}
An APD with multiplication factor $M = 10$ and excess noise factor $F = 5$ generates the same
primary photocurrent $i_\mathrm{RX}$, which is then multiplied by $M$.
The useful signal current is
\begin{equation}
  i_\mathrm{RX}^\mathrm{APD} = R_\mathrm{APD}\, P_\mathrm{RX} = R\, M\, P_\mathrm{RX} = 10 \times \SI{10}{\micro\ampere} = \SI{100}{\micro\ampere}.
\end{equation}
The avalanche multiplication is itself a random process: the APD shot noise is the PIN shot noise
amplified by $M^2$ and further increased by the excess noise factor $F$.
Over the equivalent noise bandwidth $B_\mathrm{eq} = \SI{50}{\mega\hertz}$:
\begin{equation}
\begin{aligned}
  \sigma_\mathrm{APD}^2 &= 2\, q\, M^2 F\, R\, P_\mathrm{RX}\, B_\mathrm{eq} \\
  &= 2\times\SI{1.602e-19}{\coulomb}\times 100\times 5\times\SI{e-5}{\ampere}\times\SI{5e7}{\hertz}
  = \SI{8.01e-14}{\ampere\squared}.
\end{aligned}
\end{equation}
The TIA noise is not multiplied by the avalanche gain:
$\sigma_\mathrm{TIA}^2 = (\SI{0.5}{\micro\ampere})^2 = \SI{2.5e-13}{\ampere\squared}$.
Since the two contributions are independent, their variances add:
\begin{equation}
\begin{aligned}
  \mathrm{SNR}^\mathrm{APD+TIA}_\mathrm{TOT}
  &= \frac{\left(i_\mathrm{RX}^\mathrm{APD}\right)^2}{\sigma_\mathrm{TIA}^2 + \sigma_\mathrm{APD}^2}
  = \frac{(R M P_\mathrm{RX})^2}{\sigma_\mathrm{TIA}^2 + 2 q B_\mathrm{eq} M^2 F R P_\mathrm{RX}} \\
  &= \frac{\SI{e-8}{\ampere\squared}}{\SI{2.5e-13}{\ampere\squared}+\SI{8.01e-14}{\ampere\squared}} \\
  &= \frac{\SI{e-8}{\ampere\squared}}{\SI{3.30e-13}{\ampere\squared}}
  \approx \num{3.03e4}, \\
  \mathrm{SNR}^\mathrm{APD+TIA}_\mathrm{TOT,dB} &= 10\log_{10}(\num{3.03e4}) \approx \SI{44.8}{dB}.
\end{aligned}
\end{equation}

\paragraph{Comparison.}
\begin{center}
\begin{tabular}{lccc}
  \toprule
  Receiver & Signal current & Noise variance [\si{\ampere\squared}] & SNR \\
  \midrule
  PIN + TIA & \SI{10}{\micro\ampere}  & \num{2.5e-13} (TIA) & \SI{26.0}{dB} \\
  APD + TIA & \SI{100}{\micro\ampere} & \num{2.5e-13} (TIA) $+$ \num{8.0e-14} (APD) & \SI{44.8}{dB} \\
  \bottomrule
\end{tabular}
\end{center}
The APD improves the SNR by about \SI{18.8}{dB}. If the APD added no noise, the improvement would be
$10\log_{10}(M^2) = \SI{20}{dB}$, because the signal power grows by $M^2$ while the TIA noise is unchanged.
The improvement is slightly smaller because the avalanche process adds its own shot noise,
enhanced by the factor $M^2F$, which here represents about $24\%$ of the total noise variance.
The APD is therefore advantageous when the receiver is limited by the TIA (thermal) noise.


% ---------------------------------------------------------------------
\subsection{Optical splitters}
\label{sec:2_8}
% ---------------------------------------------------------------------

\paragraph{Data.}
Transmitted power $P_\mathrm{TX} = \SI{1}{\milli\watt} = \SI{0}{dBm}$;
feeder fiber $L_1 = \SI{2}{\kilo\metre}$, drop fiber $L_2 = \SI{1}{\kilo\metre}$,
attenuation $\alpha = \SI{0.2}{dB\per\kilo\metre}$;
symmetric $1\times8$ splitter with excess loss $\varepsilon|_\mathrm{dB} = \SI{1}{dB}$.

\paragraph{Splitter model.}
A symmetric $1\times N$ splitter divides the input power equally among its $N$ outputs and
introduces an excess loss $\varepsilon$ ($\varepsilon|_\mathrm{dB} > 0$):
\begin{equation}
  P_\mathrm{out} = \frac{P_\mathrm{in}}{N}\cdot\frac{1}{\varepsilon}\ \ (\text{linear scale}),
  \qquad
  P_\mathrm{out}\big|_\mathrm{dBm} = P_\mathrm{in}\big|_\mathrm{dBm} - 10\log_{10}(N) - \varepsilon|_\mathrm{dB}.
\end{equation}
The input-output attenuation of the $1\times8$ splitter is therefore
\begin{equation}
  \mathrm{Att}|_\mathrm{dB} = 10\log_{10}(N) + \varepsilon|_\mathrm{dB}
  = 10\log_{10}(8) + 1 = 9.03 + 1 = \SI{10.03}{dB}.
\end{equation}

\paragraph{1. Received power without monitoring.}
Fiber losses: $A_1 = \alpha L_1 = \SI{0.4}{dB}$ and $A_2 = \alpha L_2 = \SI{0.2}{dB}$.

\emph{Logarithmic:} losses in dB add up:
\begin{equation}
  A_\mathrm{tot} = A_1 + \mathrm{Att}|_\mathrm{dB} + A_2 = 0.4 + 10.03 + 0.2 = \SI{10.63}{dB},
\end{equation}
\begin{equation}
  P_\mathrm{RX} = P_\mathrm{TX} - A_\mathrm{tot} = \SI{0}{dBm} - \SI{10.63}{dB} = \SI{-10.63}{dBm}.
\end{equation}

\emph{Linear:} each element multiplies the power by its linear transmission factor:
\begin{equation}
  P_\mathrm{RX} = \SI{1}{\milli\watt}\times
  \underbrace{10^{-0.04}}_{0.912}\times
  \underbrace{\tfrac{1}{8}}_{0.125}\times
  \underbrace{10^{-0.1}}_{0.794}\times
  \underbrace{10^{-0.02}}_{0.955}
  = \SI{0.0865}{\milli\watt} = \SI{86.5}{\micro\watt}.
\end{equation}

\paragraph{2. Power at the 10\% monitoring port.}
The ideal $90/10$ tap is inserted after the first fiber segment, where the power is
\begin{equation}
  P_1 = \SI{0}{dBm} - \SI{0.4}{dB} = \SI{-0.4}{dBm} = \SI{0.912}{\milli\watt}.
\end{equation}
The monitoring port receives $10\%$ of this power, i.e.\ an attenuation of
$-10\log_{10}(0.1) = \SI{10}{dB}$:
\begin{equation}
  P_\mathrm{mon} = 0.1\times\SI{0.912}{\milli\watt} = \SI{0.0912}{\milli\watt}
  \quad\Leftrightarrow\quad
  P_\mathrm{mon} = \SI{-0.4}{dBm} - \SI{10}{dB} = \SI{-10.4}{dBm}.
\end{equation}

\paragraph{3. Received power with monitoring.}
The main branch keeps $90\%$ of the power, i.e.\ an additional attenuation of
$-10\log_{10}(0.9) = \SI{0.46}{dB}$:
\begin{equation}
  A_\mathrm{tot}' = A_1 + 0.46 + \mathrm{Att}|_\mathrm{dB} + A_2 = 0.4 + 0.46 + 10.03 + 0.2 = \SI{11.09}{dB},
\end{equation}
\begin{equation}
  P_\mathrm{RX}' = \SI{0}{dBm} - \SI{11.09}{dB} = \SI{-11.09}{dBm}
  \quad\Leftrightarrow\quad
  P_\mathrm{RX}' = 0.9\times\SI{0.0865}{\milli\watt} = \SI{0.0778}{\milli\watt}.
\end{equation}

\paragraph{Summary.}
\begin{center}
\begin{tabular}{lcc}
  \toprule
  Quantity & [dBm] & [mW] \\
  \midrule
  Customer power, without tap & $-10.63$ & $0.0865$ \\
  Monitoring port (10\%)       & $-10.40$ & $0.0912$ \\
  Customer power, with tap    & $-11.09$ & $0.0778$ \\
  \bottomrule
\end{tabular}
\end{center}
Inserting the monitoring tap costs the customer only \SI{0.46}{dB}.
In exchange, the operator gets a monitoring signal of about \SI{-10.4}{dBm}, which is even stronger than the
power reaching each customer, because it is extracted before the $1\times8$ split.
